\subsection{Intersectional Subgroup Performance}
\label{sec:subgroup_performance}
\label{sec:intersectional_subgroup_performance}

\begin{table*}[htbp]
\caption{Subgroup-Level Performance for Tri-Domain Models.}
\label{tab:XI}
\centering
\begin{tabular}{llccccc}
\toprule
Classifier & Subgroup & Recall (TPR) & Precision & F1-Score & OvR Accuracy & FNR \\
\midrule
\multirow{6}{*}{\textbf{SVM (Tri)}} & White\_Males & 96.94\% & 95.36\% & 96.14\% & 98.70\% & 3.06\% \\
 & Black\_Males & 94.17\% & 95.49\% & 94.83\% & 98.29\% & 5.83\% \\
 & White\_Females & 94.72\% & 92.41\% & 93.55\% & 97.82\% & 5.28\% \\
 & Asian\_Males & 94.44\% & 92.39\% & 93.41\% & 97.78\% & 5.56\% \\
 & Asian\_Females & 92.50\% & 92.50\% & 92.50\% & 97.50\% & 7.50\% \\
 & Black\_Females & 89.44\% & 94.15\% & 91.74\% & 97.31\% & 10.56\% \\
\midrule
\multirow{6}{*}{\textbf{LR (Tri)}} & White\_Males & 96.11\% & 95.05\% & 95.58\% & 98.52\% & 3.89\% \\
 & Black\_Males & 93.06\% & 95.71\% & 94.37\% & 98.15\% & 6.94\% \\
 & Asian\_Males & 92.78\% & 90.76\% & 91.76\% & 97.22\% & 7.22\% \\
 & White\_Females & 92.22\% & 91.21\% & 91.71\% & 97.22\% & 7.78\% \\
 & Black\_Females & 91.11\% & 92.13\% & 91.62\% & 97.22\% & 8.89\% \\
 & Asian\_Females & 91.11\% & 91.62\% & 91.36\% & 97.13\% & 8.89\% \\
\bottomrule
\end{tabular}
\end{table*}

The granular classification performance profile of the best-performing tri-domain SVM model (Face $\oplus$ Emotion $\oplus$ Age) is evaluated based on the metrics in Table~\ref{tab:XI}. All subgroups achieve an F1-Score exceeding 91.00\%, ranging from 91.74\% for Black Females to 96.14\% for White Males. Nevertheless, formal fairness analysis identifies an Equal Opportunity Difference gap ($\Delta\text{TPR} = 7.50\%$) resulting from the recall disparity between White Males (96.94\%) and Black Females (89.44\% with an FNR of 10.56\%). Although the precision for Black Females remains high at 94.15\%, the lower detection sensitivity reflects intersectional phenotypic representation challenges that have not been entirely eliminated. Reporting this False Negative Rate metric balances the evaluation of OvR Accuracy (97.31\% to 98.70\%), which is prone to distortion by the 5:1 negative sample ratio dominance.

A comparison of subgroup performance profiles between the SVM model and the linear LR classifier under the tri-domain configuration is presented in Table~\ref{tab:XI}. The tri-domain LR model exhibits high performance consistency, with F1-Scores exceeding 91.00\% across all demographic classes, ranging from 91.36\% for Asian Females to 95.58\% for White Males. Formal fairness evaluation of the LR model demonstrates a narrower Equal Opportunity Difference gap ($\Delta\text{TPR} = 5.00\%$), with White Males achieving a recall of 96.11\% compared to 91.11\% for Asian Females and Black Females (FNR of 8.89\%). The White Males subgroup consistently records the highest classification performance across both classifiers. These empirical findings demonstrate that tri-domain representation fusion yields stable visual representations across all demographic subgroups in both linear and non-linear models.
